\chapter{Conclusions and future work}
\label{ch:conclusions}

This chapter summarizes the main findings of the thesis and outlines
promising directions for continuing this line of research.

\section{Conclusions}

This thesis has studied autonomous decision making in interactive
environments through the framework of games of ordered preference,
since many autonomous systems, and especially
autonomous vehicles, must satisfy multiple objectives whose relative
importance cannot be represented satisfactorily through a single
scalarized cost. Preference relations provide a suitable formalism for
capturing such prioritized objectives, but their practical value depends
on whether they can be solved efficiently, inferred from observations, and
evaluated beyond purely idealized simulation settings.

\smallskip \noindent
\textsc{computational tractability} \quad
The principal contribution of the thesis lies in the computational
tractability of games of ordered preference. In
\cref{ch:lex-ibr-ot}, we introduced lexicographic IBR over time as an
approximate solution method for receding-horizon games of ordered
preference. The results show that the method preserves the qualitative
effect of different preference relations on agent behavior while
substantially reducing compilation and solution times with respect to the
baseline formulation. This advantage becomes particularly significant as
the number of preference levels increases, which is precisely the regime
in which the standard formulation becomes least practical. The results of
this chapter therefore support the conclusion that approximate solutions
can be obtained efficiently enough to strengthen the practical viability
of preference-aware planning in interactive traffic scenarios.

\smallskip \noindent
\textsc{preference order estimation} \quad
The second contribution of the thesis concerns the inverse problem of
games of ordered preference. Although this part remains less developed
than the computational contribution, it identifies a fundamental
limitation of the complete-information assumption commonly adopted in
trajectory games. Realistic traffic scenarios do not provide direct access
to the preference order of other agents. The work initiated in
\cref{ch:estimation} therefore establishes preference-order estimation as
an essential direction for extending games of ordered preference to mixed
environments involving limited communication and human agents. In this
sense, the thesis clarifies that efficient solution methods alone are not
sufficient for practical deployment; they must be complemented by methods
capable of inferring how other agents rank their objectives.

\smallskip \noindent
\textsc{experimental validation} \quad
The third contribution of the thesis is the inclusion of experimental
validation on the Duckietown platform. The work presented in
\cref{ch:duckietown} is intentionally modest in its current scope, but it
establishes an important link between simulation and physical
experimentation. This is relevant because the significance of
preference-aware planning does not depend only on formal models and
simulation results, but also on how those methods behave under sensing
noise, actuation errors, synchronization issues, and partial
observability. The Duckietown platform should therefore be understood as
an enabling step towards a more complete sim-to-real evaluation of games
of ordered preference and preference estimation.

Taken together, the results of the thesis indicate that games of ordered
preference constitute a promising framework for representing interactive
decision making with prioritized objectives, and that their use can be
made considerably more practical through receding-horizon approximations.
At the same time, the thesis also shows that further progress in
preference estimation and experimental validation is needed if the
framework is to become fully useful in realistic autonomous-driving
settings.

\section{Future work}

Future work follows naturally from the three main axes of the thesis.
The most immediate opportunities lie in improving the tractability of
games of ordered preference further, broadening the methodological basis
for preference estimation, and strengthening the experimental validation
of the proposed framework.

\smallskip \noindent
\textsc{computational tractability} \quad
While this thesis achieved a considerable efficiency improvement by
presenting Lexicographic IBR over time, the obtained trajectories are
an approximation of the optimal solution. Future works can search for
optimizations that don't compromise accuracy. 
In particular, the tractable
reformulation proposed by Lee et al. \cite{lee:2026} replaces the
exponentially growing complete KKT system with a reduced system whose
size grows polynomially with the number of players and preference levels, all
while obtaining optimal solutions.
These results suggest that structural reformulations and temporal
approximation schemes may be combined to obtain methods that are both
efficient and accurate, while still requiring care in more general
nonlinear settings.

\smallskip \noindent
\textsc{preference order estimation} \quad
Future work should also continue the development of inverse games of
ordered preference. This includes completing the estimation methodology,
comparing alternative inference strategies, and studying how partial
observability, model mismatch, and limited sensing affect the recovery of
other agents' preference orders. Since this part of the thesis remains at
an initial stage, it offers considerable room for methodological
improvement.

\smallskip \noindent
\textsc{experimental validation} \quad
Another important direction is to make the experimental evaluation more
robust and substantially broader. This involves conducting a larger
number of experiments, considering richer interaction scenarios,
incorporating additional sensing modalities, and combining the proposed
methods with other existing works in perception, estimation, and
planning. Extending the use of the Duckietown platform in this way would
provide a more complete assessment of how games of ordered preference
behave under realistic sim-to-real conditions.
